\documentclass[reqno,11pt]{amsart}

\usepackage{amsmath,amssymb,amsthm,mathtools}
\usepackage[T1]{fontenc}
\usepackage[utf8]{inputenc}
\usepackage{booktabs}
\usepackage[colorlinks=true,linkcolor=blue,citecolor=red,urlcolor=blue]{hyperref}
\usepackage{geometry}
\geometry{a4paper,top=2.5cm,bottom=2.5cm,left=2.8cm,right=2.8cm}

\numberwithin{equation}{section}
\theoremstyle{plain}
\newtheorem{theorem}{Theorem}
\newtheorem{proposition}[theorem]{Proposition}
\newtheorem{lemma}[theorem]{Lemma}
\theoremstyle{remark}
\newtheorem{remark}[theorem]{Remark}

\newcommand{\eps}{\varepsilon}
\DeclareMathOperator{\tr}{tr}

\title[A parameter refinement for distinct zeta zeros]{A parameter refinement of the lower bound\\ for distinct zeros of the Riemann zeta function}

\author{Deep Bhattacharjee}
\address{Electro-Gravitational Space Propulsion Laboratory (EGSPL),
Bhubaneswar, Odisha, India.
ORCID: \href{https://orcid.org/0000-0003-0466-750X}{0000-0003-0466-750X}}
\email{itsdeep@live.com}

\author{Priyabrata Mandal}
\address{ORCID: \href{https://orcid.org/0000-0001-6472-6239}{0000-0001-6472-6239}}
\thanks{Corresponding author: Priyabrata Mandal.}

\date{29 September 2026. Draft, not submitted.}

\begin{document}

\begin{abstract}
Let $N(T)$ count the nontrivial zeros of $\zeta(s)$ with $0<\gamma\le T$, with
multiplicity, and let $N_d(T)$ count the distinct ones. Knausg{\aa}rd recently
proved $\liminf_{T\to\infty}N_d(T)/N(T)\ge 0.8369928814\ldots$. We show that,
with exactly the same analytic and computer-assisted inputs, the free
parameters of that argument admit a better choice, which gives
\[
\liminf_{T\to\infty}\frac{N_d(T)}{N(T)}\ \ge\ q' := \frac{62359683640669}{74504434380000}
= 0.8369929140406\ldots .
\]
The gain is about $3.26\cdot10^{-8}$. We also show that the block length
$m=1310$ used here is the largest admissible one in that argument, identify
the constraint that stops it (the residue attached to pairs of zeros off the
critical line), and record that removing this constraint alone could not
take the method beyond $0.8369964\ldots$. The result inherits every
dependency of the argument it refines, including a computer-assisted local
inequality that has not yet been refereed.
\end{abstract}

\maketitle

\noindent\textbf{MSC 2020:} 11M26 (primary); 11M06.

\section{Introduction}

It is conjectured that all nontrivial zeros $\rho=\beta+i\gamma$ of the Riemann
zeta function are simple. Let $N(T)$ count the zeros with $0<\gamma\le T$
counted with multiplicity and $N_d(T)$ the distinct ones. Montgomery's
pair-correlation method \cite{Mon73} gives lower bounds for
$N_d(T)/N(T)$. An unconditional form of Montgomery's theorem \cite{BGST24},
used as in \cite{Cla26,AF26,Lam26}, gives
$\liminf N_d/N\ge 0.83625\ldots$; Wang \cite{Wan26} improved this by about
$3\cdot10^{-8}$, and Knausg{\aa}rd \cite{Kna26} proved
\begin{equation}\label{eq:q}
\liminf_{T\to\infty}\frac{N_d(T)}{N(T)}\ \ge\ q:=\frac{16260119298029}{19426831050000}=0.83699288145\ldots
\end{equation}
by combining a mixed-multiplicity matrix inequality with a certified
seven-point local inequality from \cite{TRM26}.

The argument of \cite{Kna26} has three free parameters: a clipping
parameter $\tau\ge0$, a threshold $c$, and a block length $m$. The purpose of
this note is to observe that these can be chosen better.

\begin{theorem}\label{thm:main}
Under the inputs of \cite{Kna26} (listed in Section~\ref{sec:inputs}),
\[
\liminf_{T\to\infty}\frac{N_d(T)}{N(T)}\ \ge\ q':=\frac{62359683640669}{74504434380000}=0.83699291404068\ldots .
\]
\end{theorem}

Theorem~\ref{thm:main} improves \eqref{eq:q} by $q'-q=3.2588\ldots\cdot10^{-8}$.
The improvement is small, and we do not claim a new method. The second
result explains why the gain stops here.

\begin{proposition}\label{prop:limit}
In the argument of \cite{Kna26}, with the same window, weights and
certificate, no admissible choice of $(\tau,c,m)$ has $m>1310$.
If the admissibility condition coming from off-line pairs were dropped, the
same bound would tend to
\[
q_\infty=\frac{1+H-6/2736}{2-\delta}=0.8369964390\ldots\qquad(m\to\infty),
\]
where $H$ and $\delta$ are the constants in \eqref{eq:imports}.
\end{proposition}

\subsection*{What this does not show}
Theorem~\ref{thm:main} concerns distinct zeros, not zeros on the critical
line. More generally, a lower bound for a proportion of zeros, even one
equal to $1$, does not prove the Riemann Hypothesis: if every zero except a
set of density zero had some property, the exceptional set could still be
infinite. The statement ``$100\%$ of the zeros lie on the critical line'' in
the sense of density is compatible with infinitely many zeros off the line.
Only the statement that \emph{every} nontrivial zero lies on the line is the
Riemann Hypothesis.

\section{The argument of Knausg{\aa}rd}\label{sec:inputs}

We recall only what is needed; notation follows \cite{Kna26}. The zeros are
rescaled to mean spacing one. A fixed even probability density $f=\eta^2$ on
$[-\tfrac12,\tfrac12]$, with $K=\hat f$, gives unit vectors attached to the
zeros whose Gram matrix has entries $K(y-y')$. The operator $A$ built from
them has $\tr A=N$ and energy $E=\tr A^2$.

\subsection*{Imported inputs}
The following are used exactly as in \cite{Kna26} and are not re-proved
here:
\begin{enumerate}
\item (Energy asymptotic, \cite[Lemma 5]{BGST24}; \cite[Prop.~5]{Kna26}.)
$E_\eps=(C(f_\eps)+o(1))N$ with $C(f_\eps)\to C(f)\le 2-H$.
\item (Certified local inequality, \cite{TRM26}; \cite[Prop.~6]{Kna26}.) There
are rational weights $w_{ij}\ge0$ with $\sum_i w_{i,i+r}=2$ for $1\le r\le6$
such that for all real $y_0\le\dots\le y_6$,
\[
\tfrac{1}{2736}(y_6-y_0)+\sum_{0\le i<j\le 6}w_{ij}K(y_j-y_i)^2\ \ge\ \delta.
\]
\end{enumerate}
Here
\begin{equation}\label{eq:imports}
H=\frac{3362285207}{5000000000},\qquad \delta=\frac{891}{200000}.
\end{equation}

\subsection*{The parametrised inequality}
Let $n_1,n_2$ count the distinct on-line zeros of multiplicity $1,2$, let $h$
count distinct on-line zeros of multiplicity $\ge3$, and let $k$ count
off-line pairs, so that $N_d=n_1+n_2+h+2k$. For parameters $(\tau,c,m)$ put
\[
\Delta_m=(m-6)\delta,\qquad a=\frac{\Delta_m}{m},\qquad
\beta=\frac{6(m-6)}{2736\,m},\qquad r_h=6c-7-c^2,\qquad r_k=4c-2-c^2 .
\]
The proof of \cite[Theorem 1]{Kna26} uses the parameters only through the
following conditions:
\begin{align}
&\tau\ge0,\qquad c\ge\max(1+\tau,\,2+\tau/2), \label{eq:A1}\\
&\Delta_m\le\tau^2, \label{eq:A2}\\
&r_h-a\ge0,\qquad r_k-2a\ge0. \label{eq:A3}
\end{align}
Condition \eqref{eq:A1} is the hypothesis of the mixed-multiplicity Gram
inequality \cite[Theorem 2]{Kna26} and of the combination with the threshold
lemma \cite[Lemma 4]{Kna26}; it yields
\begin{equation}\label{eq:one}
E_\eps\ \ge\ 3N-2N_d+\tr\Psi_\tau(U)+r_hh+r_kk .
\end{equation}
Condition \eqref{eq:A2} is the hypothesis of the block dichotomy, which, with
shifted-block pinching and the window count, gives
\begin{equation}\label{eq:two}
\tr\Psi_\tau(U)\ \ge\ a\,l-\beta\,(y_l-y_1)-O_{m,\tau}(d_\eps)N-O_m(1),
\qquad l=N_d-h-2k .
\end{equation}
Combining \eqref{eq:one}, \eqref{eq:two}, the energy asymptotic and
$y_l-y_1\le N+o(N)$ gives
\begin{equation}\label{eq:final}
(2-a)N_d\ \ge\ \bigl(3-C(f_\eps)-\beta-O_{m,\tau}(d_\eps)\bigr)N+(r_h-a)h+(r_k-2a)k-o(N),
\end{equation}
and \eqref{eq:A3} lets one drop the last two terms. Letting $T\to\infty$ and
then $\eps\to0$ gives
\begin{equation}\label{eq:qm}
\liminf_{T\to\infty}\frac{N_d(T)}{N(T)}\ \ge\ q(m):=\frac{1+H-\beta}{2-a}.
\end{equation}
Knausg{\aa}rd takes $(\tau,c,m)=(12/5,17/5,1298)$, and $q(1298)=q$.

\section{Proof of Theorem~\ref{thm:main}}

Take
\[
\tau=\frac{2411}{1000},\qquad c=\frac{3411}{1000},\qquad m=1310 .
\]
Then $1+\tau=c$ and $2+\tau/2=3.2055<c$, so \eqref{eq:A1} holds.
Next,
\[
\Delta_m=\frac{145233}{25000}=5.80932<\tau^2=5.812921,
\]
so \eqref{eq:A2} holds. With
\[
a=\frac{145233}{32750000},\qquad \beta=\frac{163}{74670},
\]
one finds
\[
r_h-a=\frac{239290417}{131000000}>0,\qquad r_k-2a=\frac{5497}{26200000}>0,
\]
so \eqref{eq:A3} holds. Every step between \eqref{eq:A1}--\eqref{eq:A3} and
\eqref{eq:qm} in \cite{Kna26} is stated for general admissible parameters
(the Lean statements of \cite[\S4]{Kna26} for Theorem 2, Lemma 4 and the
counting assembly are parametric), and the smoothing error
$O_{m,\tau}(d_\eps)$ vanishes as $\eps\to0$ for any fixed $(m,\tau)$.
Therefore \eqref{eq:qm} holds with $m=1310$, and in exact arithmetic
\[
q(1310)=\frac{1+H-\beta}{2-a}=\frac{62359683640669}{74504434380000}=q'.
\]
All of these rational identities are checked by the script
\texttt{code/constants.py}, which also reproduces $q(1298)=q$. \qed

\section{Proof of Proposition~\ref{prop:limit}}

Since $\Delta_m$ increases with $m$, condition \eqref{eq:A2} forces
$\tau\ge\sqrt{\Delta_m}$. If $\tau<2$ then $\Delta_m<4$, so $m\le 903$, which
is worse than $m=1298$; hence we may assume $\tau\ge2$, where
$\max(1+\tau,2+\tau/2)=1+\tau$ and \eqref{eq:A1} reads $c\ge1+\tau$. The
function $r_k(c)=4c-2-c^2$ is decreasing for $c>2$, so the second condition
in \eqref{eq:A3} is easiest to satisfy at $c=1+\tau=1+\sqrt{\Delta_m}$, where
it becomes
\[
1+2\sqrt{\Delta_m}-\Delta_m\ \ge\ 2a=\frac{2\Delta_m}{m}.
\]
This holds for $m=1310$ and fails for every $m\ge1311$, as a direct exact
check shows (the left side decreases and the right side increases with $m$
in this range; the script \texttt{code/constants.py} performs the search).
The first condition in \eqref{eq:A3} is not binding, since
$r_h(c)>1$ for $c\in[3,3.5]$.

For the second claim, $a\to\delta$ and $\beta\to6/2736$ as $m\to\infty$, and
\[
\frac{d}{dm}q(m)>0
\]
in the relevant range, because the increase of $a$ outweighs that of
$\beta$: at leading order the sign is that of
$(1+H-\beta)\,\delta-(2-a)\cdot 6/2736>0$. Hence $q(m)$ increases to
$q_\infty$. \qed

\begin{remark}
The binding condition $r_k\ge2a$ comes from an off-line pair of
multiplicity one: in \cite{Kna26} such a pair contributes $2$ to $N_d$ and a
term with one positive eigenvalue to $Q=A-P$, which Lemma 4 of \cite{Kna26}
charges at the cost $c^2$. A sharper treatment of off-line pairs, allowing
larger $c$ (and hence larger $\tau$ and $m$), would move the bound towards
$q_\infty$. We tried absorbing a negative residue $r_k-2a<0$ using a known
lower bound for the proportion of simple zeros on the critical line (which
bounds $2k+3h\le(1-\kappa^*)N$); with either Conrey's $\kappa^*>2/5$
\cite{Con89} or the recent $\kappa^*>0.6731$ \cite{TG26}, a scan over $c\in[3.4,4.5]$ shows that
the penalty exceeds the gain whenever $r_k<2a$, so this route gives nothing.
\end{remark}

\section{Status and reproducibility}

Theorem~\ref{thm:main} is conditional on exactly the same inputs as
\cite{Kna26}: the unconditional pair-correlation asymptotic \cite{BGST24},
the interval enclosure of $C(f)$ and the certified seven-point inequality of
\cite{TRM26}, several of which are recent and not yet refereed. Nothing new
is imported. The only new content is the choice of parameters, the exact
arithmetic verifying \eqref{eq:A1}--\eqref{eq:A3}, and the optimality
analysis. The file \texttt{code/constants.py} in the accompanying repository
performs every rational computation in this note with Python's exact
\texttt{fractions} module.

\subsection*{Use of AI tools}
The parameter analysis, the verification script and a draft of this text
were produced with the assistance of an AI system (Anthropic Claude) under
the authors' direction. The authors are responsible for the content.

\begin{thebibliography}{99}

\bibitem{AF26} L.~Alp\"oge and R.~Furman, \emph{More than two thirds of the
zeta zeros are simple and on the critical line}, arXiv:2608.13637, 2026.

\bibitem{BGST24} S.~A.~C.~Baluyot, D.~A.~Goldston, A.~I.~Suriajaya and
C.~L.~Turnage-Butterbaugh, \emph{An unconditional Montgomery theorem for pair
correlation of zeros of the Riemann zeta-function}, Acta Arith.\ \textbf{214}
(2024), 357--376.

\bibitem{Cla26} Claude (Anthropic), \emph{More than two thirds of the zeros of
the Riemann zeta function lie on the critical line}, preprint, 2026.

\bibitem{Con89} J.~B.~Conrey, \emph{More than two fifths of the zeros of the
Riemann zeta function are on the critical line}, J.\ Reine Angew.\ Math.\
\textbf{399} (1989), 1--26.

\bibitem{Kna26} K.~M.~Knausg{\aa}rd, \emph{More than 83.69\% of the zeros of
the Riemann zeta function are distinct}, arXiv:2609.33043, 2026.

\bibitem{Lam26} Y.~Lamzouri, \emph{A new proof that more than $2/3$ of the
zeros of the Riemann zeta function are simple and on the critical line},
arXiv:2609.02882, 2026.

\bibitem{Mon73} H.~L.~Montgomery, \emph{The pair correlation of zeros of the
zeta function}, Proc.\ Sympos.\ Pure Math.\ \textbf{24}, AMS, 1973, 181--193.

\bibitem{TG26} tawanerguo-cn/zeta-simple-zeros, \emph{A certified
unconditional 67.3192911473\% lower bound for simple zeros of the Riemann
zeta function}, research repository, 2026,
\url{https://github.com/tawanerguo-cn/zeta-simple-zeros}.

\bibitem{TRM26} trmdy/zeta-simple-zeros-673137, commit 1610b97, 2026,
\url{https://github.com/trmdy/zeta-simple-zeros-673137}.

\bibitem{Wan26} B.~Wang, \emph{Proportions of the non-trivial zeros of the
Riemann zeta function}, arXiv:2609.24167, 2026.

\end{thebibliography}

\end{document}
